
\begin{answerBox}{2in}          % please don't edit this line, doing so will throw off the page layout

  %%%%%%%%%%%%%%%%%%%%%%%%%%%%%%%%%%%%%%%%%%%%%%%%%%%%%%
  %%%%%%%%%%%%%%%%%%%%%%%%%%%%%%%%%%%%%%%%%%%%%%%%%%%%%%
  %%
  %% PROBLEM 4 COLLABORATORS AND CITATIONS 
  %%
  %%%%%%%%%%%%%%%%%%%%%%%%%%%%%%%%%%%%%%%%%%%%%%%%%%%%%%
  %%%%%%%%%%%%%%%%%%%%%%%%%%%%%%%%%%%%%%%%%%%%%%%%%%%%%%
  
  \paragraph{Problem \arabic{problem} collaboration and citations}~\GradescopeComment
  {This should be on page 8 (bottom) of your PDF.}

  I collaborated on this problem with the following people:
  \begin{blue}
    \begin{itemize}
    \item \REPLACEME            % write "none"if none
    \end{itemize}
  \end{blue}

  My answers use ideas from the following sources (other than the lecture notes and textbooks): 
  \begin{blue}
    \begin{itemize}
    \item \REPLACEME            % write "none"if none
    \end{itemize}
  \end{blue}

\end{answerBox}

\newpage

\paragraph{Problem 4 answer}\GradescopeComment{This should be on page 9 (top) of your PDF.}

\begin{enumerate}[label=(\alph*)]

\ITEM 
  %%%%%%%%%%%%%%%%%%%%%%%%%%%%%%%%%%%%%%%%%%%%%%%%%%%%%% 
  %%%%%%%%%%%%%%%%%%%%%%%%%%%%%%%%%%%%%%%%%%%%%%%%%%%%%% 
  %% 
  %% PROBLEM 4(a) ANSWER 
  %% 
  %%%%%%%%%%%%%%%%%%%%%%%%%%%%%%%%%%%%%%%%%%%%%%%%%%%%%% 
  %%%%%%%%%%%%%%%%%%%%%%%%%%%%%%%%%%%%%%%%%%%%%%%%%%%%%% 

  % Find a maximum-value $s$-$t$ flow $f$ in the network.
  % Draw the network, labeling each edge $(u, w)$ with the flow $f(u,w)$ on the edge.

  Here is the flow:

\begin{answerBox}{3in}          % please don't edit this line, doing so will throw off the page layout

  \begin{center}
    \begin{tikzpicture}[
        vertex/.style={circle, inner sep=0pt, minimum width=15pt, draw}, 
        label/.style={midway, fill=white, inner sep=2pt}, 
        every path/.style={->}
        ]
      \node[vertex] (s) {s}; 
      \node[vertex] (b) [right =of s] {b}; 
      \node[vertex] (a) [above =of b] {a}; 
      \node[vertex] (c) [below =of b] {c}; 
      \node[vertex] (d) [right =of a, xshift=20pt] {d}; 
      \node[vertex] (f) [right =of b, xshift=20pt] {f}; 
      \node[vertex] (g) [right =of c, xshift=20pt] {g}; 
      \node[vertex] (t) [right =of f] {t}; 

      %% EDGES AND FLOWS START HERE:

      % replace each question mark below by the flow on the edge:
      \path (s) edge node [label] { ? } (a);  % replace "?" with flow on edge s->a
      \path (s) edge node [label] { ? } (b);  % replace "?" with flow on edge s->b
      \path (s) edge node [label] { ? } (c);  % replace "?" with flow on edge s->c

      \path (a) edge node [label] { ? } (b);  % etc...
      \path (a) edge node [label] { ? } (d);  
      \path (a) edge node [label] { ? } (f); 

      \path (b) edge node [label] { ? } (f) ;
      \path (b) edge node [label] { ? } (c) ; 

      \path (c) edge node [label] { ? } (g) ; 

      \path (d) edge node [label] { ? } (t) ; 
      \path (d) edge node [label] { ? } (f) ; 

      \path (f) edge node [label] { ? } (t); 
      \path (f) edge node [label] { ? } (g) ; 

      \path (g) edge node [label] { ? } (t) ; 
      \path (g) edge node [label] { ? } (b) ; 
    \end{tikzpicture}

    % Or you can delete the block above, and instead use some other method to draw the network,
    % get it into a PDF, then include the PDF in your LaTeX document using these commands:

    % \begin{center}
    %   \includegraphics[height=3in]{YOUR_FILE_NAME_HERE}  % (omit the .pdf extension in the name)
    % \end{center}

    % See https://texblog.org/2012/02/23/crop-figures-with-includegraphics
    % if you want to crop the included image.

  \end{center}
\end{answerBox}


\item
  %%%%%%%%%%%%%%%%%%%%%%%%%%%%%%%%%%%%%%%%%%%%%%%%%%%%%% 
  %%%%%%%%%%%%%%%%%%%%%%%%%%%%%%%%%%%%%%%%%%%%%%%%%%%%%% 
  %% 
  %% PROBLEM 4(b) ANSWER 
  %% 
  %%%%%%%%%%%%%%%%%%%%%%%%%%%%%%%%%%%%%%%%%%%%%%%%%%%%%% 
  %%%%%%%%%%%%%%%%%%%%%%%%%%%%%%%%%%%%%%%%%%%%%%%%%%%%%% 

  % What is the value of the flow? 

  The value of the flow is
  \BLUE{\[
      \REPLACEME.
    \]}

\item

  %%%%%%%%%%%%%%%%%%%%%%%%%%%%%%%%%%%%%%%%%%%%%%%%%%%%%% 
  %%%%%%%%%%%%%%%%%%%%%%%%%%%%%%%%%%%%%%%%%%%%%%%%%%%%%% 
  %% 
  %% PROBLEM 4(c) ANSWER 
  %% 
  %%%%%%%%%%%%%%%%%%%%%%%%%%%%%%%%%%%%%%%%%%%%%%%%%%%%%% 
  %%%%%%%%%%%%%%%%%%%%%%%%%%%%%%%%%%%%%%%%%%%%%%%%%%%%%% 

  % [review 2026-09-27] fix: N: comment had the complement backwards: 'T=S\setminus V' -> 'T=V\setminus S'
  % Find an $s$-$t$ cut $(S, T=V\setminus S)$ of minimum capacity in the network.
  % For the cut you found, what is $S$?

  The $s$-$t$ cut is $(S, T = V\setminus S)$ where
  \[
    S = \{ s,
    \BLUE{
      \REPLACEME
    }
    \}.
  \]


\item
  %%%%%%%%%%%%%%%%%%%%%%%%%%%%%%%%%%%%%%%%%%%%%%%%%%%%%% 
  %%%%%%%%%%%%%%%%%%%%%%%%%%%%%%%%%%%%%%%%%%%%%%%%%%%%%% 
  %% 
  %% PROBLEM 4(d) ANSWER 
  %% 
  %%%%%%%%%%%%%%%%%%%%%%%%%%%%%%%%%%%%%%%%%%%%%%%%%%%%%% 
  %%%%%%%%%%%%%%%%%%%%%%%%%%%%%%%%%%%%%%%%%%%%%%%%%%%%%% 

  % And what is the capacity of the cut?

  The capacity of the cut is \BLUE{\[
      \REPLACEME.
      \]}
  
\end{enumerate}


%%  THE COLLABORATORS AND CITATIONS STATEMENT IS AT THE TOP OF THIS FILE


%%% Local Variables:
%%% mode: latex
%%% TeX-master: "homework_9"
%%% End:
